% !PRINTED 695
\weberpage{695}
%<<CONT
enclosed, and in absolute value
\[
R_m < \frac{2\,C}{m^s},
\]
or, if $s > c$,
\[
R_m < \frac{2\,C}{m^c}.
\]

This upper limit for $R_m$, which is independent of $s$, can however, if $c$ is positive, be pressed down below any value however small by taking $m$ sufficiently large.

But from this there follows not only the convergence, but also the continuity of $S$. For if one takes $m$ sufficiently large, then not only $R_m$, but also the fluctuation of $R_m$ under varying $s$ becomes infinitely small, and $S_m$ is, for a fixed $m$, a continuous function of $s$. Hence $S$ too is continuous.
% [sic] The German source has a singular subject (die Schwankung) with the plural verb werden.

It is also to be remarked here that the assumption about $\sigma_m$ by no means presupposes the convergence of the infinite series $\Sigma a_k$. It is not even necessary that the $a_k$ be real; for if they are imaginary, one need only apply the theorem to the real and to the imaginary part. Finally, it is also not necessary to assume the $a_1, a_2, a_3, \dots$ to be constants. Everything remains valid if they are \gesperrt{continuous functions} of $s$.

2. \gesperrt{If $\mu_1, \mu_2, \dots, \mu_n, \dots$ denotes an infinite set of positive numbers of such a nature that two finite numbers $\alpha, \beta$ can be assigned in such a way that for every $n$, however large,}
\weq{3}{\alpha < \dfrac{n}{\mu_n} < \beta,}
\gesperrt{then the infinite series}
\[
S = \frac{1}{\mu_1^s} + \frac{1}{\mu_2^s} + \frac{1}{\mu_3^s} + \cdots
\]
\gesperrt{is convergent so long as the exponent $s$ is greater than 1, and the product $(s-1)\,S$ likewise remains between the limits $\alpha$ and $\beta$ when $s$ approaches the value 1 sufficiently closely, or at least approaches one of these values without limit.}
